\documentclass[11pt,a4paper]{article}
\usepackage[utf8]{inputenc}
\usepackage[T1]{fontenc}
\usepackage{lmodern}
\usepackage{amsmath,amssymb,amsthm,mathtools}
\usepackage{geometry,booktabs,longtable,array,enumitem,microtype,seqsplit}
\usepackage{xcolor}
\usepackage{hyperref}
\geometry{margin=2.65cm}
\setlength{\emergencystretch}{2em}
\newcommand{\ttbr}[1]{\texttt{\seqsplit{#1}}}
\newcommand{\N}{\mathbb N}
\newcommand{\Z}{\mathbb Z}
\newcommand{\R}{\mathbb R}
\newcommand{\Q}{\mathbb Q}
\newcommand{\U}{\mathcal U}
\newtheorem{theorem}{Theorem}[section]
\newtheorem{proposition}[theorem]{Proposition}
\newtheorem{lemma}[theorem]{Lemma}
\newtheorem{corollary}[theorem]{Corollary}
\theoremstyle{definition}
\newtheorem{definition}[theorem]{Definition}
\newtheorem{remark}[theorem]{Remark}
\hypersetup{colorlinks=true,linkcolor=blue!45!black,citecolor=blue!45!black,
 urlcolor=blue!45!black,pdftitle={Finite Orbit Budgets, Arithmetic Transport Barriers, and Paired-Path Certificates for Collatz},pdfauthor={Piero Borgatta}}
\title{Finite Orbit Budgets, Arithmetic Transport Barriers,\\
and Paired-Path Certificates for Collatz\\[0.4em]
\large Version 7: a technical continuation of the phantom-shadowing programme}
\author{Piero Borgatta\thanks{Independent researcher.
Correspondence: \texttt{info@pieroborgatta.com}.
ORCID: \href{https://orcid.org/0009-0001-8025-2405}{0009-0001-8025-2405}.}}
\date{16 November 2026}
\begin{document}
\maketitle

\begin{abstract}
This technical continuation of version~6 reports finite arithmetic results
and a more precise research boundary for a Lean-assisted Collatz programme.
For every finite injective shortcut orbit we prove an explicit reciprocal
budget smaller than~81. It yields a uniform first-arrival correction product
at most $e^{27}$ and exact height constraints on existing inverse-word
transports. A quadratic occupation inequality handles repeated sources and
cyclic targets, and gives finiteness of every fixed bounded-slope catalogue;
its time bound is existential, not an effective search cutoff.
We distinguish these results from conditional mass arguments. A single
lower logarithmic-density input pays finite prefix-free off-path catalogues,
including a possible entry into a cycle, without identifying a canonical
harmonic trace. On paper, two classical paired-word rules are composed
into an explicit mixed language with a nearby smaller companion and an
attributed terminal merger. A minimum divergent unit seed must leave that
language. Several exits admit exact local charges against subsequent
forced contractions, but their global composition remains open.
The manuscript supplies proofs, attribution, a comparison with version~6,
and a source-bound verification record. The finite core is checked in
the pinned Lean~4.29.1 development; the paired-path and conditional results
are paper proofs. No exclusion of all divergent orbits or all nontrivial
cycles, universal smaller-source cover, or solution of Collatz is claimed.
\end{abstract}

\begingroup\small\tableofcontents\endgroup

\section{Purpose, contribution, and status}
\label{sec:scope}
Version~6 \cite{borgatta_v6} established the local exact-cylinder and
precision-budget interfaces used here, corrected an earlier global
continuity claim, and explained why an unrestricted infinite-code
exclusion would restate the divergence problem. This version is a
standalone continuation, not a replacement of the frozen version~6 files.
It selects results from the subsequent development that have explicit
mathematical statements and an inspectable proof boundary.

There are three different kinds of assertion in this paper.
\begin{enumerate}[leftmargin=*]
\item The finite estimates in Section~\ref{sec:v7-finite-core} have
corresponding checked declarations in the local Lean project. Their paper
proofs make the estimates readable without requiring Lean.
\item Sections~\ref{sec:transport-barriers}, \ref{sec:conditional-budget},
and~\ref{sec:paired-paths} give paper proofs. They are not represented as
new checked Lean theorems. The logarithmic-density premise in
Section~\ref{sec:conditional-budget} is stated as an external input.
\item The final research obligations are questions. A finite certificate
interface, a locally correct paired rule, or a positive result on a
constructed family is not a proof of coverage of every integer orbit.
\end{enumerate}

The formalization and AI-assisted review are methods of producing and
checking statements, not grounds for a claim of mathematical priority.
Parity cylinders, inverse trees, multiplicative potentials, and paired
parity words have substantial precedents
\cite{terras1976,lagarias2003,krasikovlagarias2003,garner1985,eliatucker2015}.
Qualitative integer Green finiteness was already explicit in
\cite[Proposition~11]{weaver9v1}. We therefore claim a local formal
integration, explicit finite constants, and the specific proved
interfaces below; we do not claim the first general use of these tools.

\subsection{What changed after version~6}
\begin{center}\small
\begin{tabular}{>{\raggedright\arraybackslash}p{0.22\linewidth}>{\raggedright\arraybackslash}p{0.33\linewidth}>{\raggedright\arraybackslash}p{0.34\linewidth}}
\toprule
Topic & Version~6 or earlier boundary & Present continuation\\
\midrule
Exact words & Final odd bit and exact residue cylinder & Retained as the legality test for all recodings\\
Precision & One fixed expansive word consumes precision; unrestricted switches can reset it & First-arrival product and explicit finite transport constraints\\
Orbit packing & No pointwise descent cover & Universal finite injective reciprocal budget; no new descent cover\\
Inverse catalogues & Arithmetic interfaces and selected searches & Finite bounded-slope catalogue, including returns; no effective time bound\\
Mass comparison & Relative-mass lower bound missing & Smaller conditional dependency interface; lower construction still missing\\
Paired paths & No global rank for arbitrary switching & Mixed exact rules, terminal merger, necessary exits, local contraction charges\\
\bottomrule
\end{tabular}
\end{center}

\subsection{Conventions}
Natural numbers used as orbit states are positive. The shortcut map is
$T(n)=n/2$ for even $n$ and $(3n+1)/2$ for odd $n$; the Syracuse map
is $S(n)=(3n+1)/2^{a(n)}$ on positive odd integers, with
$a(n)=\nu_2(3n+1)\ge1$. A \emph{unit} is a positive odd integer
not divisible by~3. Write $\U$ for this set. Every positive odd
Syracuse image is a unit, even when its source is divisible by~3.
An orbit point is \emph{periodic} if it has a positive return to itself.
A nonperiodic point may still be preperiodic; this distinction matters
in the finite argument allowing a later entry into a cycle.
A divergent orbit has no repeated state and tends to infinity.

The gain of a forward word of length $k$ and exponent sum $A$ is
$\Omega=3^k/2^A$. A gain is a homogeneous affine slope, not the
actual ratio of the terminal and initial states. The latter includes
the positive correction product. Inverse words are read from target
to seed; prefix-free always refers to that orientation.

\section{Finite orbit budgets and exact predecessor catalogues}
\label{sec:v7-finite-core}

This section records finite arithmetic statements in the post-v6 Lean
development. Their proofs do not import a predecessor-density theorem,
an atom floor, a mixing estimate, or a conjecture about divergent orbits.
The parity-cylinder geometry is classical
\cite{terras1976,lagarias2003}. Uniform reciprocal and correction-product
bounds for finite distinct trajectories are already proved in
\cite[Section~2]{valbuena2026}. A related separation by odd-step count
appears in Curry's collision-free window estimate
\cite[Theorem~2.3]{curry2026}. The contribution asserted here is the
specific parity-moment proof, its explicit rational constant, and its
integration into the local formal development. Qualitative Green
finiteness is not claimed as new.

Write
\[
 T(n)=\begin{cases}n/2,&n\text{ even},\\(3n+1)/2,&n\text{ odd},\end{cases}
 \qquad
 S(n)=\frac{3n+1}{2^{\nu_2(3n+1)}}\quad(n>0\text{ odd}).
\]
For an actual Syracuse segment $x_i=S^i(u)$, put
\[
 A_k=\sum_{i<k}\nu_2(3x_i+1),\qquad
 \Omega_k=\frac{3^k}{2^{A_k}}.
\]
For an exact exponent word $w$, the corresponding seed-independent gain
is $\Omega(w)=3^{|w|}/2^{\sum w}$. An exact word includes the prescribed
valuations and the odd terminal state, rather than only an affine identity.
A first arrival at $n$ at time $k$ means $x_k=n$ and $x_i\ne n$ for
every $i<k$. The odd-unit restriction $3\nmid u$ used in the subsequent
research programme is not needed for most estimates below; where only
positive odd seeds are specified, multiples of three are allowed as seeds.

\subsection{An explicit reciprocal budget for finite injective orbits}

\begin{theorem}[Finite reciprocal budget]\label{thm:v7-reciprocal}
Let $u>0$ and $L\ge0$. If $T^0(u),\ldots,T^{L-1}(u)$ are distinct, then
\begin{equation}\label{eq:v7-reciprocal}
 \sum_{i=0}^{L-1}\frac1{T^i(u)}
 \le\frac{133547}{1655}<81.
\end{equation}
No assumption is made about the orbit after time $L-1$.
\end{theorem}
\begin{proof}
Let $J_N(s)$ count the odd entries among
$T^0(s),\ldots,T^{N-1}(s)$. The exact parity lift and its moment identity
are
\begin{align}
 T^N(s+2^N)&=T^N(s)+3^{J_N(s)},
 &J_N(s+2^N)&=J_N(s),\label{eq:v7-parity-lift}\\
 \sum_{0\le s<2^N}q^{J_N(s)}&=(1+q)^N.\label{eq:v7-parity-moment}
\end{align}
The first identity follows by induction through the $N$ exposed parity
bits. The last newly exposed bit is complemented by adding $2^N$;
pairing the two lifts therefore multiplies the preceding moment by
$1+q$, proving the second identity. Here the auxiliary representative
$s=0$ is permitted, with $T(0)=0$.
Induction also gives
\[
 2^N T^N(s)<3^{J_N(s)}(s+2^N),
 \qquad 0\le s<2^N,
\]
so $T^N(s)<2\cdot3^{J_N(s)}$.

Consider a finite set $B$ in $[2^N,2^{N+1})$ on which $T^N$ is
injective. Write $n=2^N+s$. At a fixed value $J_N(s)=m$, distinct $n$
give distinct reduced endpoints $T^N(s)$ by
\eqref{eq:v7-parity-lift}. There are at most $2\cdot3^m$ such
endpoints. Consequently the number of elements with $J_N(s)<M$ is
less than $3^M$. On the other hand, substituting $q=3/2$ in
\eqref{eq:v7-parity-moment} bounds the number with $J_N(s)\ge M$ by
$(5/2)^N(2/3)^M$. Thus
\[
 \frac{\#B}{2^N}\le
 \frac{3^M}{2^N}+\left(\frac54\right)^N\left(\frac23\right)^M.
\]

Let $b_N$ count the states of the segment in the $N$th dyadic band.
Delete its last $N$ time indices, or the whole segment if it has fewer
than $N$ entries. The map $T^N$ is injective on the remaining states:
their images are still distinct states of the original segment.
Restoring the deleted indices costs at most $N$. With
$M=\lfloor3N/5\rfloor$, it follows that
\begin{align}
 \frac{b_N}{2^N}&\le\frac{N}{2^N}+h_N+\ell_N,
 \label{eq:v7-band-bound}\\
 h_N&=\left(\frac54\right)^N\left(\frac23\right)^{\lfloor3N/5\rfloor},
 \qquad \ell_N=\frac{3^{\lfloor3N/5\rfloor}}{2^N}.\notag
\end{align}
Only finitely many bands are occupied. Their reciprocal contribution is
at most $\sum_N b_N/2^N$. Grouping the upper envelopes into blocks of
five gives
\[
 h_{5r+s}=\left(\frac{3125}{3456}\right)^r h_s,
 \qquad
 \ell_{5r+s}=\left(\frac{27}{32}\right)^r\ell_s
 \quad(0\le s<5),
\]
with $\sum_{s<5}h_s=3271/576$ and
$\sum_{s<5}\ell_s=51/16$. Finite geometric sums therefore give the
respective bounds $19626/331$ and $102/5$. The terminal error costs at
most $\sum_N N/2^N\le2$. At $N=0$, however, the actual band contains
only the possible state $1$, whereas \eqref{eq:v7-band-bound} assigns
two units. Saving that extra unit yields
\[
 \sum_{i<L}\frac1{T^i(u)}
 \le 2+\frac{19626}{331}+\frac{102}{5}-1
 =\frac{133547}{1655}.
\]
All estimates used in the implementation are finite rational inequalities;
the displayed geometric bounds do not require an asymptotic counting
hypothesis.
\end{proof}

The principal declaration is
\ttbr{CollatzShadowing.OrbitReciprocal.finite\_orbit\_reciprocal\_sum\_le}.
The companion real-scalar theorem is
\ttbr{finite\_orbit\_reciprocal\_sum\_real\_lt\_eighty\_one}
in the same namespace. The temporal deletion, which is essential for the
injectivity assertion used above, is formalized in
\texttt{FiniteOrbitTransfer.lean}; injectivity of the full Collatz map on
arbitrary integer sets is not assumed.

\subsection{First arrivals and the additive correction product}

\begin{proposition}[First-arrival envelope]\label{prop:v7-first-arrival}
For every positive odd first-arrival segment $u\to n$ of length $k$,
\begin{equation}\label{eq:v7-product}
 n=\Omega_k u \mathcal P_k,
 \qquad
 \mathcal P_k=\prod_{i<k}\left(1+\frac1{3x_i}\right),
 \qquad 1\le \mathcal P_k\le P_0:=\exp(27).
\end{equation}
In particular, $\Omega_k u\le n\le P_0\Omega_k u$.
\end{proposition}
\begin{proof}
A finite trajectory through its first arrival at a target is injective,
including its endpoint. Indeed, a repetition before the endpoint would
make the remaining future repeat and would produce an earlier visit to
that same target. Every Syracuse step at an odd state has positive
exponent, and the elapsed shortcut time at $S^i(u)$ is exactly $A_i$.
All strictly interior shortcut states of such an episode are even.
Hence a first Syracuse arrival at the odd target $n$ is a first shortcut
arrival at time $A_k$. Its shortcut segment is injective, and the
Syracuse states $x_i$, $i<k$, are a subset of that segment. Theorem
\ref{thm:v7-reciprocal} gives $\sum_{i<k}1/x_i<81$.

Multiplying the exact identities
$x_{i+1}=(3/2^{a_i})x_i(1+1/(3x_i))$ proves
\eqref{eq:v7-product}. Each factor is at least one, and
$1+t\le e^t$ gives
$\mathcal P_k\le\exp(\sum_{i<k}1/(3x_i))\le e^{27}$.
\end{proof}

The first-arrival sum and exact shortcut-clock interface are the
declarations
\ttbr{FirstArrivalBudget.first\_arrival\_reciprocal\_sum\_lt\_eighty\_one}
and \ttbr{FirstArrivalBudget.shortcut\_clock}, under the common
\texttt{CollatzShadowing} prefix. The product identity and its exponential
estimate are
\ttbr{WordProductEnvelope.inflation\_eq\_product} and
\ttbr{WordProductEnvelope.inflation\_le\_exp\_of\_reciprocal\_budget}.
The identity itself allows repeated visits; the uniform constant $P_0$
here is justified by the first-arrival hypothesis. It must not be
reused for arbitrarily many turns around a cycle.

\begin{corollary}[Height of an existing recoding]
\label{cor:v7-height}
Suppose $w:u\to n$ is as in Proposition~\ref{prop:v7-first-arrival}
and $w':v\to y$ is another exact positive odd-source word. If
$\Omega(w')\ge\Omega(w)/C$ for $C>0$, then
\begin{equation}\label{eq:v7-height}
 v\le C\frac{y}{n}P_0u.
\end{equation}
In particular, equal word lengths and equal exponent totals allow $C=1$.
If additionally $y\le21n$, then $v\le21e^{27}u$.
\end{corollary}
\begin{proof}
The nonnegative affine constant of the exact word $w'$ gives
$\Omega(w')v\le y$. The first-arrival product for $w$ gives
$1/\Omega(w)=u\mathcal P_k/n\le P_0u/n$. Combining the two inequalities proves
\eqref{eq:v7-height}.
\end{proof}

This is implemented by
\ttbr{WordProductEnvelope.comparable\_slope\_height} and, for the
same-profile specialization,
\ttbr{FirstArrivalBudget.first\_arrival\_same\_profile\_height}.
The canonical unit sibling of a unit $n$ is $4n+1$ if it is a unit,
and $16n+5$ otherwise; both possibilities are at most $21n$.
The canonical-sibling specialization is
\ttbr{FirstArrivalBudget.canonical\_first\_arrival\_profile\_height}.
These are bounds on a recoding \emph{that exists}. They supply neither
that recoding nor a positive proportion of successfully recoded sources.

\subsection{Squared occupation, including repeated cyclic visits}

Let $\phi_k=\Omega_k/x_k$. The exact decrement of its square is
\begin{equation}\label{eq:v7-square-loss}
 \phi_k^2-\phi_{k+1}^2
 =\Omega_k^2\frac{6x_k+1}{x_k^2(3x_k+1)^2}.
\end{equation}
This follows directly from
$\phi_{k+1}=3\Omega_k/(3x_k+1)$ and does not require first arrival.
Put
\[
 V(n)=\frac{n^2(3n+1)^2}{6n+1}\quad(n>0).
\]

\begin{theorem}[Finite quadratic occupation]\label{thm:v7-square}
For a fixed positive odd seed $u$, a target $n>0$, and any finite set
$H$ of times with $S^k(u)=n$ for all $k\in H$,
\begin{equation}\label{eq:v7-visit-bound}
 \sum_{k\in H}\Omega_k^2\le\frac{V(n)}{u^2}.
\end{equation}
For any finite catalogue $\mathcal W$ of distinct exact words ending at
$n$, with arbitrary positive odd sources, even when sources repeat,
\begin{equation}\label{eq:v7-catalogue-square}
 \sum_{w\in\mathcal W}\Omega(w)^2\le2V(n).
\end{equation}
Periodic targets and the empty word are allowed.
\end{theorem}
\begin{proof}
The nonnegative losses \eqref{eq:v7-square-loss} telescope from
$\phi_0^2=1/u^2$. At a visit to $n$, the loss is exactly
$\Omega_k^2/V(n)$, proving \eqref{eq:v7-visit-bound} for every finite
set of visiting times. From one fixed source there is only one actual
word of any specified length; distinct words in its fibre therefore
correspond to distinct visiting times, including returns. Group the
catalogue by source and use
\[
 \sum_{u\in U}\frac1{u^2}\le2
\]
for any finite set $U$ of positive integers. The latter estimate follows,
for example, by telescoping $1/u^2\le2/u-2/(u+1)$.
\end{proof}

The exact loss, repeated-visit bound, and arbitrary-source catalogue
bound are respectively
\ttbr{QuadraticOccupation.potential\_sq\_sub\_succ},
\ttbr{QuadraticOccupation.sum\_weight\_sq\_visits\_le}, and
\ttbr{QuadraticOccupation.tilted\_catalogue\_le}.
The simpler fixed-depth estimate $\sum\Omega(w)^2\le2n^2$ is
\ttbr{IntegerAtomTilt.tilted\_integer\_atom\_le}; it uses source
injectivity at that depth. The same constant is shared across arbitrary
depths if the target has no positive return, as stated by
\ttbr{IntegerAtomTilt.tilted\_nonreturn\_catalogue\_le}.
At a cyclic target this additional hypothesis cannot be dropped from
the $2n^2$ argument; Theorem~\ref{thm:v7-square} supplies the separate
repeated-visit estimate.

The corresponding infinite integer Green sum was already shown finite
in \cite[Proposition~11]{weaver9v1}. Equation
\eqref{eq:v7-catalogue-square} is a finite elementary realization with
an explicit constant; it is not a priority claim for Green finiteness.

\begin{corollary}[Finite bounded-slope catalogue]
\label{cor:v7-bounded-catalogue}
For every positive integer target $n$ and positive rational $R$, the
set of all exact words from positive odd sources to $n$ satisfying
$\Omega(w)\ge1/R$ is finite, and
\begin{equation}\label{eq:v7-catalogue-card}
 \#\mathcal W_R(n)\le2V(n)R^2.
\end{equation}
Every corresponding source satisfies $u\le Rn$. There exist integers
$K,A$ bounding all their word lengths and exponent totals, respectively.
The conclusion includes cyclic targets.
\end{corollary}
\begin{proof}
Each selected word contributes at least $1/R^2$ to
\eqref{eq:v7-catalogue-square}. Hence every finite subcatalogue has
cardinality at most the fixed number $2V(n)R^2$; an infinite set would
have finite subsets of arbitrarily large cardinality, a contradiction.
The source bound follows from $\Omega(w)u\le n$. The length and total
bounds are maxima of two functions on the resulting finite set.
\end{proof}

The five public results in \texttt{BoundedSlopeCatalogue.lean} are
\ttbr{subcatalogue\_card\_le}, \ttbr{catalogue\_finite},
\ttbr{catalogue\_ncard\_le}, \ttbr{catalogue\_parameter\_bounds}, and
\ttbr{catalogue\_source\_bound}. Its implemented catalogue is slightly
larger than the odd-source catalogue displayed here: only positivity of
the source is imposed. The numerical source and cardinality bounds are
explicit. The proof supplies the temporal bounds $K,A$ only
\emph{existentially}; neither a numerical hitting-time cutoff nor a
procedure whose stopping criterion follows from the cardinality bound
is asserted.

\subsection{Finite horizons and the direction of collision certificates}

Two additional interfaces delimit what these finite upper bounds can
and cannot provide. First, every positive odd segment obeys
\begin{equation}\label{eq:v7-shifted-growth}
 2^k(S^k(u)+1)\le3^k(u+1).
\end{equation}
Thus a word of length at most $K$ ending at $n$ has source at least
$R\ge1$ whenever $3^K(R+1)\le2^K(n+1)$. For a finite catalogue with
distinct sources at least $R$,
\[
 \sum_w\Omega(w)^2\le n^2\sum_u\frac1{u^2}
 \le\frac{2n^2}{R}.
\]
Distinctness is automatic at one fixed depth, and across arbitrary
depths at a target without positive return. These are
\ttbr{FiniteHorizonBudget.iterate\_shifted\_growth},
\ttbr{tilted\_tail\_catalogue}, \ttbr{tilted\_fixed\_depth\_tail},
and \ttbr{tilted\_nonreturn\_horizon}. Repeated sources at a cyclic
target require their own occupation budget; they are not silently
treated as distinct.

Second, let $R_0,\ldots,R_N$ be nondecreasing real numbers with
$R_i\ge a>0$. Suppose a finite subprobability catalogue $P_i$ contains
a distinguished branch with mass $R_i/R_{i+1}$, and let
$D_i=1-\sum_{j\in P_i}p_{i,j}^2$. Then
\begin{equation}\label{eq:v7-collision-budget}
 0\le D_i\le1-\left(\frac{R_i}{R_{i+1}}\right)^2,
 \qquad
 \sum_{i<N}D_i\le\frac2a(R_N-R_0).
\end{equation}
Indeed $1-r^2\le2(1-r)$, and
$(R_{i+1}-R_i)/R_{i+1}\le(R_{i+1}-R_i)/a$; the increments telescope.
The finite algebra is
\ttbr{SpineDispersion.collision\_spine\_budget}.

If two distinct actual branches each have mass at least $c$, every
branch has mass at most $1-c$, whence $D_i\ge c$. This is
\ttbr{SpineDispersion.collision\_defect\_ge\_of\_two\_branches};
the more general lower-allocation certificate is
\ttbr{collision\_defect\_ge\_lower\_allocations}. These statements
do not prove that such branches exist. Moreover, a truncated catalogue
has \emph{larger} defect than an extension, because adding squared
masses decreases $1-\sum p^2$. Its computed defect is therefore an
\emph{upper}, not a lower, bound on the full law's defect. A branch
lower certificate transfers only when it remains valid with the actual
unnormalized masses for every finite extension and the corresponding
limit is justified. Renormalized sample frequencies do not supply this
certificate.

The named declarations are tied to the repository's pinned verification
record. The seven-module integration on 15 November audited all 136
compiled declarations in that scope; the subsequent four-module
integration audited all 68 declarations, including five encoded-private
helpers, and passed an incremental full-library build of 3404 jobs.
The permitted transitive axioms in those audits were
\texttt{propext}, \texttt{Classical.choice}, and \texttt{Quot.sound}.
A separate Ubuntu replay reproduced the latter audit and finite
arithmetic outputs. Those recorded checks use Lean and Mathlib
\texttt{v4.29.1}; they are not a cold rebuild of all dependencies or a
replay with an independently implemented kernel. The release manifest
identifies the sources and evidence included in this version.


\section{Arithmetic transport barriers and the tower boundary}
\label{sec:transport-barriers}

\subsection{A small new seed requires an early change}
\begin{proposition}[A preserved-prefix obstruction]
Suppose two distinct positive odd seeds $b,v$ realize the same first
$r$ Syracuse exponents, of total $A$. Then
\[
 2^{A+1}\mid(v-b).
\]
Consequently a witness with $v\le C_b j\log j$ can preserve at most
$O_b(\log j)$ initial exponent bits. The conclusion is unaffected
by allowing a different final word length.
\end{proposition}
\begin{proof}
The affine constants cancel, giving
$S^r(v)-S^r(b)=3^r(v-b)/2^A$. Both endpoints are odd, so the
difference is even; $3$ is invertible modulo every power of~$2$.
Thus $2^{A+1}\le |v-b|\le b+C_b j\log j$.
\end{proof}
This is the classical cylinder constraint in a useful orientation, not
a new cylinder theorem. It rules out changing only a short terminal
window while retaining a long forward prefix, if a genuinely small
new seed is required. It does not rule out a global change beginning
near the source. If the rival word ends earlier, the shared prefix is
only defined up to that endpoint.

\subsection{Uniform routing does not buy a contracting return}
Uniform legal modular routing with linear total valuation is available
in \cite[Theorem~8]{weaver48}; its absolute constants are not numerically tabulated.
The following elementary count concerns the real cost of a particular
way of using such a primitive.
\begin{proposition}[Fixed-total routing cost]\label{prop:routing-cost}
Fix length $b\ge1$, total $E\ge b$, precision $r\ge1$, and a unit
residue $a$ modulo $3^r$. Suppose words of length $b$ and total $E$
route every positive odd unit root into that residue. Then
\[
 \binom{E-1}{b-1}\ge2\cdot3^{b+r-1},\qquad
 \frac{2^E}{3^b}\ge4\cdot3^{r-1}.
\]
If routing into $-1$ modulo $3^{k+1}$ is followed by $k$ inverse
steps of exponent~$1$, the composed inverse real slope is at least
$4\cdot2^k$.
\end{proposition}
\begin{proof}
A word has the inverse form $I_w(x)=(2^Ex-C_w')/3^b$.
Prescribing its terminal residue modulo $3^r$ fixes one root residue
modulo $3^{b+r}$, because $2^E$ is invertible. There are
$2\cdot3^{b+r-1}$ unit root residues and only
$\binom{E-1}{b-1}$ positive compositions of $E$ into $b$ parts.
Legality restrictions cannot increase coverage.
Use $\binom{E-1}{b-1}\le2^{E-1}$. Finally,
$I_1^k(y)+1=(2/3)^k(y+1)$; the extra terminal ternary digit
preserves the unit guard and gives the stated slope.
\end{proof}
For fixed parameters the menu is finite. Its offsets cannot overcome
this slope at all arbitrarily large roots. This is not a pointwise
lower bound on every routed seed and not a no-go for arbitrary
root-dependent clocks. If only a fraction $\delta$ of the root
residues must be covered, the same count has the additional factor
$\delta$, leaving a possible role for rare conditioning.

\subsection{What exact tower information does and does not say}
The post-v6 tower is
\[
 Q_v(q)=\frac{9^{2^v q}-1}{2^{v+3}},\qquad q>0\text{ odd},
 \qquad Q_v=Q_v(1).
\]
For distinct positive odd parameters $q,r$, the elementary lifting
identity gives
\[
 \nu_2\bigl(Q_v(q)-Q_v(r)\bigr)=\nu_2(q-r).
\]
Indeed factor out $9^{2^v\min(q,r)}$ and use
$\nu_2(9^m-1)=3+\nu_2(m)$ for nonzero integer $m$.
This exact isometry, its finite residue bijections, and finite-word
transport are integrated in \texttt{TowerParameter.lean} and
\texttt{TowerSuffix.lean}. They explain why a finite suffix can be
realized by choosing its parameter. They do not imply that the one
fixed parameter $q=1$ realizes an arbitrary suffix, or that one
positive parameter realizes an arbitrary infinite word.

The source-specific squaring and carry constraints retain more
information than a generic cylinder. The archived post-v6 experiments
nevertheless refuted several proposed finite completion filters, not
the tower itself. None supplies a uniform bound for the total parity
weight of the exact tower. The release preserves these negative
results as research history; it does not turn another exact source
description into a new dynamical inequality.

\section{A conditional finite budget for off-path certificates}
\label{sec:conditional-budget}
This section is a paper argument. Its external analytic input is
separated from the unconditional first-arrival estimate. It avoids
assuming natural-density existence or identifying a canonical trace.

\begin{definition}[Logarithmic-density input]\label{def:log-input}
For a nonperiodic unit $v$, let $\mathcal B(v)$ be the set of positive
odd integers whose forward Syracuse orbit reaches $v$, allowing
zero time. Assume there is $\alpha>0$ such that for every such $v$,
\begin{equation}\label{eq:log-input}
 \liminf_{X\to\infty}\frac{2}{\log X}
 \sum_{\substack{u\le X\\u\in\mathcal B(v)}}\frac1u
 \ge\frac{\alpha}{4v}.
\end{equation}
\end{definition}
The lower logarithmic statement is already given from an atom-floor
input in \cite[Section~2]{weaver12}. The corresponding atom-floor
manuscript is \cite{eljaouhari2026}. These are external preprint
inputs here, not theorems imported into the local Lean development.
Their published formal artifacts were not rebuilt for this release.
All subsequent uses of~\eqref{eq:log-input} remain conditional on it.

\begin{lemma}[Finite antichain packing]\label{lem:log-antichain}
Under~\eqref{eq:log-input}, a finite forward antichain $V$ of
nonperiodic units satisfies $\sum_{v\in V}1/v\le4/\alpha$.
\end{lemma}
\begin{proof}
If two predecessor basins intersect, determinism makes their targets
forward comparable. Thus these basins are disjoint. The normalized
sum of reciprocals of all positive odd integers tends to~$1$.
Apply finite additivity and the lower bounds in~\eqref{eq:log-input};
no interchange of an infinite family and a liminf is used.
\end{proof}

\begin{theorem}[A budget with the cycle status undecided]
\label{thm:finite-offpath}
Let $n_0=b,n_1,\ldots,n_J$ be an actual Syracuse segment of units
with distinct states. For each $1\le j\le J$, take a finite family
$\mathcal W_j$ of nonempty actual words from unit seeds $v_w$ to
$n_j$, satisfying all three conditions:
\begin{enumerate}[label=(\roman*),leftmargin=*]
\item each word is a first arrival at $n_j$;
\item its last parent differs from the observed parent $n_{j-1}$;
\item the family is prefix-free when read inversely from $n_j$.
\end{enumerate}
Write $\Omega_j$ for the observed gain and
$\Gamma_w=\Omega(w)/\Omega_j$. Under~\eqref{eq:log-input},
\begin{equation}\label{eq:finite-budget}
 \sum_{j=1}^J\sum_{w\in\mathcal W_j}\Gamma_w
 \le b e^{27}\left(\frac8\alpha+1\right).
\end{equation}
If all observed targets are nonperiodic, $4be^{27}/\alpha$ suffices.
No decision about the orbit after the displayed prefix is required
as an input to the finite certificate.
\end{theorem}
\begin{proof}
For both first-arrival segments,
$n_j=b\Omega_jP_j=v_w\Omega(w)P_w$, with
$1\le P_j,P_w\le e^{27}$ by
Proposition~\ref{prop:v7-first-arrival}. Hence
$\Gamma_w=bP_j/(v_wP_w)\le be^{27}/v_w$.

If the targets are all nonperiodic, so are the seeds. At a common
target, prefix-freeness and first arrival make the seeds a forward
antichain. Between two distinct off-path ingress times, a seed
already entering the observed path must thereafter use the observed
parent at the later target. Therefore such seeds are incomparable,
and Lemma~\ref{lem:log-antichain} pays the whole family once.

If the prefix contains periodic points, let $i$ be the first such
index. For $j>i$, the observed parent is the unique cyclic parent
of $n_j$. An off-path word must enter the cycle for the first time
at $n_j$ itself; otherwise its last parent would be cyclic. For
$j<i$ the targets and seeds are nonperiodic. All seeds for
$j\ne i$ form one nonperiodic antichain: an earlier path enters
the cycle at $n_i$, whereas a later off-path ingress enters at a
different displayed cycle point. Distinctness of the observed
segment is used here. If $i=0$, no exceptional catalogue is present.

For $i\ge1$, separate $\mathcal W_i$. Its nonperiodic seeds form
a second antichain, by prefix-freeness and first arrival. A seed
from this family may also occur in a later catalogue, which is why
one cannot simply pay all terms once. Periodic seeds in
$\mathcal W_i$ lie on the unique inverse cyclic chain; prefix-freeness
allows at most one of them, with reciprocal at most~$1$.
The two nonperiodic budgets and this one possible term give
$\sum 1/v_w\le8/\alpha+1$, proving~\eqref{eq:finite-budget}.
\end{proof}

For a hypothetical divergent path, a construction of such catalogues
whose total relative gain is unbounded would contradict the theorem.
This is a sufficient research criterion, not an already proved lower
bound. With a rival word of length $\ell$ and total $B$, its cost is
\[
 \Gamma=3^{\ell-j}2^{A_j-B}=2^{-E},\qquad
 E=(B-A_j)-(\ell-j)\log_2 3.
\]
The clock is free. Moreover $\Gamma=bP_j/(vP_w)$, so the cost is
the logarithm of the new seed within a bounded correction. Renaming
that seed as a debt does not make its construction easier.

\section{Paired paths and the cost of leaving a block language}
\label{sec:paired-paths}

This section gives paper proofs, not new Lean declarations. Paired parity
words and their concatenation are classical tools
\cite{garner1985,eliatucker2015}. Our purpose is to give a small, explicit
instance with positive integer states, unit repairs, a different final
parent, and quantified exit costs. Neither the mechanism nor the first
construction of a smaller source before descent is claimed as new;
the latter already has substantially stronger constructions in
\cite[Theorem~3]{weaver21}. The open issue here is coverage of an already
specified integer orbit.

\subsection{An exact block relation}

For a positive-exponent word $w=(a_1,\ldots,a_k)$, write
$A(w)=\sum_i a_i$ and
\[
 F_w(x)=\frac{3^k x+C_w}{2^{A(w)}},\qquad
 C_{()}=0,\qquad C_{w\mathbin{\|}(a)}=3C_w+2^{A(w)}.
\]
The function $F_w$ is an affine expression; it equals $S^k$ only on the
exact valuation cylinder of $w$. We use the exact-cylinder criterion
established in v6 \cite{borgatta_v6}: for a positive odd integer $x$,
the word $w$ is realized precisely when $F_w(x)$ is an odd integer.
For completeness, the reverse implication follows by induction from
$C_w=3^{k-1}+2^{a_1}C_{(a_2,\ldots,a_k)}$: oddness of the final quotient
forces the first valuation to be $a_1$, after which the same argument
applies to the remaining word. All actual intermediate states are
positive because the initial state is positive.

\begin{lemma}[Translation of exact blocks]\label{pp:translation}
Let $W,V$ have the same length $k\ge1$ and total exponent $A$.
For an even integer $c$,
\[
 F_V(x+c)=F_W(x)+c
 \quad\Longleftrightarrow\quad
 C_V-C_W=c(2^A-3^k).
\]
If this identity holds, $x$ realizes $W$, and both $x+c$ and
$F_W(x)+c$ are positive, then $x+c$ realizes $V$ exactly.
\end{lemma}
\begin{proof}
Compare the affine constants. The second assertion follows because
$F_W(x)+c$ is then a positive odd integer and the exact-cylinder
criterion applies to the positive odd input $x+c$.
\end{proof}

The translation $H(x)=x-2$ gives two rules:
\begin{equation}\label{pp:rules}
 W_1=(1,2)\longmapsto V_1=(2,1),\qquad
 W_2=(2,1,2)\longmapsto V_2=(1,1,3).
\end{equation}
Their affine constants and cylinders are
\[
\begin{array}{c|c|c|c|c}
 W&V&(3^k,2^A)&(C_W,C_V)&\text{source cylinders }W,V\\ \hline
 (1,2)&(2,1)&(9,8)&(5,7)&11,9\pmod {16}\\
 (2,1,2)&(1,1,3)&(27,32)&(29,19)&57,55\pmod {64}.
\end{array}
\]
In each row $F_V(x-2)=F_W(x)-2$. The least positive inputs
in these cylinders exceed $2$, and the block endpoints also exceed $2$.
Consequently any actual concatenation $p\in\{W_1,W_2\}^*$ starting
at $b>2$ is followed by an actual concatenation $\sigma(p)$ starting
at $b-2$, with a difference of $2$ at every block boundary.
The substitution $\sigma$ is defined by \eqref{pp:rules}.
This is intertwining between different affine maps. It does not require
one companion map to commute with every observed block. Indeed, the
fixed points of $F_{W_1}$ and $F_{W_2}$ are $-5$ and $29/5$.
The second block has gain $27/32<1$; the entire block language is not
an expansive language.

\subsection{An attributed terminal word and a mixed family}

The binary word used below comes from Cassian Mahler's manuscript
\cite[Section~4, Equation~(4.4)]{mahler152}. Its broader search and
survival certificates are not inputs to the following direct proof.
Set
\begin{align*}
 w_*={}&(1,1,2,1,2,1,2,1,2,1,2,1,2,2,2,2,1,2,1,2,2,1),\\
 g_*={}&(4,\underbrace{1,\ldots,1}_{15},2,2,1,4,3,3).
\end{align*}
Both have length $22$ and total exponent $34$. The recursive definition
of the affine constant gives
\[
 C_{w_*}=118309227953,\qquad C_{g_*}=181071347171,
 \qquad C_{g_*}-C_{w_*}=2\cdot3^{22}.
\]
Thus $F_{g_*}(x-2)=F_{w_*}(x)$. These integers provide a direct finite
check of the displayed identity; they do not certify the other results
of \cite{mahler152}.

\begin{proposition}[A mixed paired family]\label{pp:mixed-family}
Suppose a positive odd integer $b$ realizes $p w_*$, where
$p\in\{W_1,W_2\}^*$. Then $b-2$ realizes $\sigma(p)g_*$ and the
two paths have the same endpoint, length, total exponent, and gain.
Their last parents are different. For a unit input $b$, a unit source
for the alternative path is given by the following table; $\Gamma$
denotes its gain divided by the observed gain:
\[
\begin{array}{c|c|c}
\text{condition}&\text{alternative unit source}&\Gamma\\ \hline
b\equiv1\pmod3&b-2&1\\
b\equiv2\pmod3,\ p\text{ begins with }W_1&(3b-5)/4&4/3\\
b\equiv2\pmod3,\ p\text{ begins with }W_2&(3b-5)/2&2/3\\
b\equiv2\pmod3,\ p=()&(3b-5)/16&16/3.
\end{array}
\]
The last three paths have one fewer odd step.
\end{proposition}
\begin{proof}
The seed cannot be $1$: its valuation word consists of $2$'s, whereas
$p w_*$ contains a $1$. Hence $b>2$, so the companion seed is positive.
Iterate Lemma~\ref{pp:translation} along $p$, then use the displayed
terminal identity. Exact oddness of its endpoint proves the legality
of $g_*$. The observed last exponent is $1$, whereas the alternative
last exponent is $3$. Equal last parents would have the same valuation,
so the parents are different.

If $b\equiv1\pmod3$, the source $b-2$ is a unit. If
$b\equiv2\pmod3$, it is a positive odd multiple of $3$.
Delete its first actual step. That step has exponent $f=2,1,4$
in the three respective cases, and its endpoint is
$(3b-5)/2^f$. This endpoint is a positive odd unit. Deleting the
step multiplies the gain by $2^f/3$ and leaves the last parent
unchanged. The remaining length is at least $21$.
\end{proof}

The third unit source in the table can exceed $b$. It is not a
unit descent certificate. A useful subclass avoids this issue. Put
\[
 A=W_1^4,\qquad B=W_1^3 W_2.
\]
These blocks form a prefix code, have lengths $8,9$ and exponent totals
$12,14$, and are not powers of one common word. Every nonempty prefix
of either block has gain greater than $1$. For the last three letters
of $B$, the smallest intermediate gain is
$(9/8)^3(3/4)=2187/2048>1$; its total gain is $19683/16384>1$.
The exponent totals of the prefixes of $w_*$ are
\[
 (1,2,4,5,7,8,10,11,13,14,16,17,19,21,23,25,26,28,29,31,33,34).
\]
Each is below the corresponding $i\log_2 3$; their smallest gain is
$3^{21}/2^{33}>1$. Hence every actual word in
$\{A,B\}^*w_*$ stays above its initial source at every nonempty
prefix, while Proposition~\ref{pp:mixed-family} supplies a smaller
unit source with the same future. This includes the empty macro prefix.

These are identities on proper cylinders, not coverage of a supplied
infinite orbit. The relative odd Haar measure of all cylinders
$\{W_1,W_2\}^*w_*$ is only
\[
 \sum_{r\ge0}(2^{-3}+2^{-5})^r2^{-34}
 =\frac{2^{-29}}{27}.
\]
They are disjoint because $w_*$ starts with $11$, whereas the two
blocks start with $12$ or $212$. Nor does a merger alone certify a
first arrival. The first-arrival and prefix-free guards must be checked
when the earlier mass budgets are applied; for a nonperiodic target,
an actual arrival is automatically its unique arrival from that source.

\subsection{A minimum divergent seed must leave the language}

\begin{proposition}[A necessary exit; no analytic input]
\label{pp:minimum-exit}
If a divergent positive odd unit orbit exists, let $b$ be its least
unit starting value. Its infinite valuation word does not belong to
$\{W_1,W_2\}^{\mathbb N}$.
\end{proposition}
\begin{proof}
The seed satisfies $b>1$. Its first valuation must be $1$, since a
valuation at least $2$ gives the smaller divergent unit
$S(b)\le(3b+1)/4<b$. If its word were an infinite block concatenation,
its first block would therefore be $W_1$.

At block boundaries the actual orbit of $b-2$ differs from that of $b$
by $2$. Those values tend to infinity. An orbit of a deterministic map
on the positive integers that has an unbounded subsequence cannot
repeat a state, and therefore also tends to infinity. Thus $b-2$
diverges. If it is a unit, this contradicts minimality. Otherwise its
first valuation is the first alternative valuation of $V_1$, namely
$2$, and $S(b-2)=(3b-5)/4$ is a divergent unit smaller than $b$.
This is again a contradiction.
\end{proof}

The proposition requires neither a merger nor an atom-floor theorem.
It constrains the \emph{minimum} divergent unit seed, not every
divergent seed. It leaves open all exits not covered by a further
certificate. Elia and Tucker's counterexamples to a specific proposed
classification of paired stems \cite{eliatucker2015} are a relevant
warning: a collection of valid local rules does not establish their
completeness.

\subsection{Local costs and subsequent forced contraction}

At a block boundary write the pair as $(x,x-2)$, and let $x'$ be the
next observed state. If the observed exponent is $e\ge2$, the
companion exponent is exactly $1$ and the new relation is
$y'=2^{e-1}x'-3$. If $e=1$, the companion exponent is
$f=1+\nu_2(x'-3)\ge2$ and $y'=(x'-3)/2^{f-1}$.
These follow directly by subtracting $6$ from $3x+1$.
After the prefix $21$ the relation is $y=2x-5$.
An exit $13$ or $213$ consequently produces
\[
 y'=\operatorname{oddpart}(z-1),\qquad z=x' >1.
\]
Here $z=1$ is incompatible with an integral preceding state of
valuation $3$ in either case.

\begin{lemma}[A forced run]\label{pp:forced-two}
If $z=2^t u+1$, where $u$ is positive odd and $t\ge1$, then the next
$r=\lfloor(t-1)/2\rfloor$ valuations are all $2$ and
\[
 S^r(z)=(3/4)^r(z-1)+1.
\]
\end{lemma}
\begin{proof}
For $s\ge3$,
$3(2^s u+1)+1=4(3\cdot2^{s-2}u+1)$, with the factor in parentheses
odd. Iterate while the remaining exponent of $2$ is at least $3$.
\end{proof}

\begin{proposition}[Local cost or a drop below the reference seed]
\label{pp:buffer-cost}
Let $b>1$ be a reference value and suppose an observed path reaches
$z=2^t u+1$. Either one of $z,S(z),\ldots,S^r(z)$ lies below $b$,
where $r$ is as in Lemma~\ref{pp:forced-two}, or
\begin{equation}\label{pp:t-charge}
 t\le 2+\frac{2\log((z-1)/(b-1))}{\log(4/3)}.
\end{equation}
In particular the inequality holds along the orbit of a minimum
divergent unit seed $b$ at an exit $13$ or $213$.
For an exit beginning with $11$, if the first companion valuation is
$f\ge7$, the corresponding bound is
\begin{equation}\label{pp:f-charge}
 f\le 8+\frac{2\log((z-1)/(b-1))}{\log(4/3)}
\end{equation}
unless the observed path subsequently falls below $b$.
\end{proposition}
\begin{proof}
If there is no drop, the last forced state gives
$(3/4)^r(z-1)+1\ge b$. Hence
$r\log(4/3)\le\log((z-1)/(b-1))$, and $t\le2r+2$ proves
\eqref{pp:t-charge}. All later states of a minimum divergent unit
orbit are units at least $b$.

For the last assertion put $u=S(x-2)$, where $x$ is the boundary
state. From $3(x-2)+1=2^f u$, the first observed states are
\[
 x_1=3+2^{f-1}u,\quad x_2=5+3\cdot2^{f-2}u,\quad
 z=x_3=1+9\cdot2^{f-6}u.
\]
For $f\ge7$ the three observed exponents are exactly $1,1,4$.
Apply \eqref{pp:t-charge} with $t=f-6$.
\end{proof}

The cost in these inequalities is controlled by the actual height
buffer $\log((z-1)/(b-1))$. It is not an $O(\log j)$ cost at time $j$.
The remaining cases $f=3,4,5,6$ are four parametric classes, with
unbounded integer states, not four finite numerical states.

\begin{proposition}[A local loss of the same companion]
\label{pp:collapse-family}
For every fixed $K\ge5$, there are arbitrarily large positive odd
units $b$ whose observed word begins with $W_1^K(1,3)$, all its
nonempty prefixes have gain greater than $1$, and the companion
$b-2$ reaches $1$ while the observed path reaches an arbitrarily
large value.
\end{proposition}
\begin{proof}
Choose arbitrarily large $t=6m$ satisfying
$3^{2K+3}\mid 2^{t+1}+7$. Such $m$ exist because $64$ has order
$3^{N-2}$ modulo $3^N$ on $1+9\mathbb Z/3^N\mathbb Z$, for $N\ge2$,
as follows from $\nu_3(64^h-1)=2+\nu_3(h)$. The target $-7/2$
belongs to that subgroup. Put
\[
 x=\frac{2^{t+4}+11}{9},\qquad
 b=\frac{8^{K+1}(2^{t+1}+7)}{9^{K+1}}-5.
\]
For sufficiently large $t$ these are positive odd integers, with
$b\equiv1\pmod3$, so both initial seeds are units. The identity
$x+5=(9/8)^K(b+5)$ and the exact block cylinders give
\[
 b\xrightarrow{W_1^K}x\xrightarrow{(1,3)}2^t+1,
 \qquad
 b-2\xrightarrow{V_1^K}x-2\xrightarrow{(2,t+2)}1.
\]
The last intermediate states are $(2^{t+3}+7)/3$ and
$(2^{t+2}-1)/3$, which verify the last valuations exactly.
The smallest gain on the final two observed letters is
$(9/8)^K(9/16)>1$ for $K\ge5$, since $9^6>16\cdot8^5$.
All preceding prefixes are also expansive.
\end{proof}

This is not a divergent orbit construction or a counterexample to
Collatz. For fixed $K$, the ratio $(2^t)/(b-1)$ stays bounded as
$t\to\infty$, while the following forced run has unbounded length.
It therefore eventually drops below $b$. The example rules out a
uniformly bounded local cost for keeping this particular companion;
it does not rule out paying the cost using later contraction.

The next required result is a composition theorem: at every exit,
either certify a drop below the original seed, or choose another
positive unit companion and control the accumulated costs. The local
bounds above do not prove that result, since later expansive blocks
can replenish the height buffer. No universal descent, no complete
exit classification, and no nonsummable allocation of alternative
paths follows from this section.


\section{Literature boundary and the next mathematical obligation}
\label{sec:next}
The rStation references are version-pinned manuscripts attributed to
their named platform accounts. Their platform acceptance is not treated
as evidence of independent peer review or formal correctness. Results
used only as precedents are distinguished from the explicit external
premise of Section~\ref{sec:conditional-budget}.

Tao's theorem \cite{tao2022} gives almost-bounded minima in logarithmic
density, not descent for every positive integer. Predecessor-count
and canonical-trace results likewise have their own domains and
quantifiers. In particular, fixed-root nonatomicity of a canonical
inverse law is already available from ancestor decay
\cite{valbuena2026,weaver12}; it does not supply a uniform critical
clock for moving targets. The nonperiodic linear trace upper bound
must not be extended to arbitrary cycles without its return factor.

Several tempting subproblems are already solved in the literature.
Chosen absorbing words of linear length are given in
\cite{weaver81}, and modular routes with linear total valuation in
\cite{weaver48}. An absorbing word selects a compatible source
cylinder; it does not prescribe the future bits of a supplied
integer. A modular route controls a residue, not the height of the
new seed. Proposition~\ref{prop:routing-cost} explains one specific
failure of combining those tools without conditioning.

A recent manuscript \cite{weaver150} makes a possible divergent
orbit's exceptional ternary behaviour quantitative: any weak limit
of its empirical Syracuse-state measures has a positive singular
component, with the stated lower fraction about~0.08855. This is a
necessary restriction, not a divergence exclusion. None of our
pointwise floor, telescope, or finite packing estimates forces the
opposite regularity. We do not identify these empirical measures
with an inverse Doob law or a Haar-stationary control model.

The concrete continuation suggested by Section~\ref{sec:paired-paths}
is to handle exits from the paired language by one of three verified
outcomes: a merger, a controlled change to another companion, or
descent below the original minimum seed. The large-parameter exits
already display a charge against later forced contractions. Those
charges depend on a growth buffer, not on $O(\log j)$, and expansive
blocks can replenish that buffer. A global amortization has therefore
not been proved. Adding more terminal words without such an argument
would repeat the incomplete-classification problem illustrated by
\cite{eliatucker2015}.

The cycle problem is separate. Finite product constraints and
restricted-height cycle exclusions do not exclude every positive
cycle. No analytic result cited here is promoted into the checked
library merely by being referenced, and no large external
computational certificate is claimed as independently reproduced.

\section{Verification, release contents, and limitations}
\label{sec:verification}
The release binds the paper, its English supporting audits, the
canonical Lean sources, build pins, selected strict dependency
audits, and their recorded outputs in a reproducible source archive.
The canonical package has~90 Lean files including its entrypoint;
26 files are generated finite-certificate modules. The release
manifest, rather than a mutable branch name, identifies the bytes.
The standalone \texttt{CollatzShadowing} repository is a documented
export of this package; it is not an independent proof implementation.

The fresh pinned incremental build completes~3404 jobs. It reuses
the existing dependency cache and is not a cold rebuild of Mathlib.
The selected new arithmetic declarations are checked transitively
using only \texttt{propext}, \texttt{Classical.choice}, and
\texttt{Quot.sound}. Historical receipts cover the earlier post-v6
modules; their source hashes are checked individually, since later
imports changed the main entrypoint. The whole library additionally
contains a documented finite list of earlier \texttt{native\_decide}
proofs. Their native compilation assumptions are not part of the
strict arithmetic-core claim.

Lean and Mathlib are pinned to~4.29.1. This predates the documented
kernel soundness fix in~4.32.2 and further soundness fixes and defensive
checks in~4.34.0 \cite{lean4322,lean4340}. No exploit is established
in this development. The reported result is acceptance by the pinned
toolchain, with explicit dependency audits, not an independent
current-kernel validation or a claim that pinning removes those
limitations. Migration and a fresh independent checker replay remain
separate verification work. Source scans for placeholders alone do
not settle this trust question.

AI assistants supported literature retrieval, drafting, symbolic
derivations, Lean development, and adversarial review. Several agents
reviewed one another's work in shared context; they are not independent
human referees. The named human author directs the programme and is
responsible for the manuscript. Cross-review and successful finite
checks do not certify mathematical novelty. Authorship and the wider
methodology are described in the repository's \texttt{METHODOLOGY.md}.

\section{Conclusion}
The finite reciprocal and occupation budgets make the arithmetic
interfaces more explicit and reproducible. The conditional antichain
argument reduces a sufficient mass criterion to one stated external
input, while the paired-word construction demonstrates a legal
source-specific mechanism on a proper language. Its exits already
show why local correctness and a global descent argument are different
tasks. The missing results are a coverage theorem and a global cost
bound, or another argument that genuinely excludes the exceptional
integer orbit. They remain open.

\begingroup\small
\begin{thebibliography}{99}
\raggedright

\bibitem{borgatta_v6}
P.~Borgatta, \emph{Exact Exponent Cylinders, Precision Budgets, and Singular
Coding in Accelerated Syracuse Dynamics}, version 6.0.0, Zenodo preprint,
24 September 2026.
\href{https://doi.org/10.5281/zenodo.22936057}{doi:10.5281/zenodo.22936057}.

\bibitem{terras1976}
R.~Terras, \emph{A stopping time problem on the positive integers},
Acta Arithmetica \textbf{30} (1976), no. 3, 241--252.
\href{https://doi.org/10.4064/aa-30-3-241-252}{doi:10.4064/aa-30-3-241-252}.

\bibitem{lagarias2003}
J.~C.~Lagarias, \emph{The $3x+1$ problem: An annotated bibliography
(1963--1999) (sorted by author)}, arXiv:math/0309224,
first submitted 13 September 2003; version 13, 11 January 2011.
\href{https://arxiv.org/abs/math/0309224v13}{Pinned arXiv version}.

\bibitem{krasikovlagarias2003}
I.~Krasikov and J.~C.~Lagarias, \emph{Bounds for the $3x+1$ problem using
difference inequalities}, Acta Arithmetica \textbf{109} (2003), no. 3,
237--258.
\href{https://doi.org/10.4064/aa109-3-4}{doi:10.4064/aa109-3-4};
\href{https://arxiv.org/abs/math/0205002v1}{arXiv:math/0205002v1}.

\bibitem{garner1985}
L.~E.~Garner, \emph{On heights in the Collatz $3n+1$ problem},
Discrete Mathematics \textbf{55} (1985), no. 1, 57--64.
\href{https://doi.org/10.1016/S0012-365X(85)80020-0}{doi:10.1016/S0012-365X(85)80020-0}.

\bibitem{eliatucker2015}
M.~Elia and A.~Tucker, \emph{Consecutive integers and the Collatz conjecture},
Integers \textbf{15} (2015), Article A54, 30 December 2015.
\href{https://math.colgate.edu/~integers/p54/p54.pdf}{Journal article};
\href{https://arxiv.org/abs/1511.09141v1}{arXiv:1511.09141v1},
30 November 2015.

\bibitem{curry2026}
M.~J.~Curry, \emph{An Explicit Windowed Sparsity Bound for Divergent
$3x+1$ Orbits, with Logarithmic-Floor Exclusion beyond the Harmonic
Barrier}, Zenodo preprint, 24 August 2026.
\href{https://doi.org/10.5281/zenodo.22087163}{doi:10.5281/zenodo.22087163}.

\bibitem{valbuena2026}
L.~Valbuena, \emph{Collatz matrix obstructions, synchronization, and basin
densities: three papers and Lean artifacts}, software and manuscript
archive, version 2026.09.28.5, Zenodo, 28 September 2026.
\href{https://doi.org/10.5281/zenodo.23025885}{doi:10.5281/zenodo.23025885}.
The manuscript \emph{Collatz basin densities and canonical traces from
first-passage stabilization} is cited at
\href{https://github.com/x1xhlol/collatz-two-coordinate-obstruction/blob/a7275d5a73814d3d9a400c6fe6577de9542b03a3/paper/tao-basin-densities.tex}{source snapshot \texttt{a7275d5a}}.

\bibitem{weaver9v1}
Lyapunov Weaver, \emph{Ternary Syracuse laws: tail conditioning, moment
growth, and predecessor bounds}, rStation Collatz paper 9, version 1,
21 October 2026, 04:48 UTC, preprint.
\href{https://rstation.ai/problems/collatz/papers/9/v1}{Version page};
\href{https://rstation.ai/api/v1/papers/collatz_9/versions/1/download}{pinned source}.

\bibitem{weaver12}
Lyapunov Weaver, \emph{First-landing bounds and canonical inverse laws
for Collatz}, rStation Collatz paper 12, version 1,
21 October 2026, 08:59 UTC, preprint.
\href{https://rstation.ai/problems/collatz/papers/12/v1}{Version page};
\href{https://rstation.ai/api/v1/papers/collatz_12/versions/1/download}{pinned source}.

\bibitem{weaver21}
Lyapunov Weaver, \emph{Smaller Collatz sources before first descent:
power gains and conditional exponent shells}, rStation Collatz paper 21,
version 3, 23 October 2026, 19:46 UTC, preprint.
\href{https://rstation.ai/problems/collatz/papers/21/v3}{Version page};
\href{https://rstation.ai/api/v1/papers/collatz_21/versions/3/download}{pinned source}.

\bibitem{weaver48}
Lyapunov Weaver and Fourier Wren, \emph{Hierarchical Syracuse renewal
and linear valuation budgets for modular routing}, rStation Collatz
paper 48, version 2, 28 October 2026, 05:46 UTC, preprint.
\href{https://rstation.ai/problems/collatz/papers/48/v2}{Version page};
\href{https://rstation.ai/api/v1/papers/collatz_48/versions/2/download}{pinned source}.

\bibitem{weaver81}
Lyapunov Weaver, \emph{Linear absorbing words and protected
smaller-source certificates for Collatz}, rStation Collatz paper 81,
version 2, 1 November 2026, 21:59 UTC, preprint.
\href{https://rstation.ai/problems/collatz/papers/81/v2}{Version page};
\href{https://rstation.ai/api/v1/papers/collatz_81/versions/2/download}{pinned source}.

\bibitem{weaver150}
Lyapunov Weaver, \emph{Critical barriers and memory costs in Collatz
inverse dynamics}, rStation Collatz paper 150, version 2,
15 November 2026, 21:53 UTC, preprint.
\href{https://rstation.ai/problems/collatz/papers/150/v2}{Version page};
\href{https://rstation.ai/api/v1/papers/collatz_150/versions/2/download}{pinned source}.

\bibitem{mahler152}
Cassian Mahler, \emph{Slope--Ballot Duality, Exact Binary Escape
Thresholds, and Parity-Trap Shifts Exceeding the $3/4$ Collatz Survival
Barrier}, rStation Collatz paper 152, version 1,
15 November 2026, 21:47 UTC, preprint.
\href{https://rstation.ai/problems/collatz/papers/152/v1}{Version page};
\href{https://rstation.ai/api/v1/papers/collatz_152/versions/1/download}{pinned source}.

\bibitem{eljaouhari2026}
N.~El Jaouhari, \emph{Uniform Syracuse atoms and exact density regularity
with quantitative Collatz predecessor counts}, version 1, Zenodo
preprint, 14 October 2026.
\href{https://doi.org/10.5281/zenodo.23377671}{doi:10.5281/zenodo.23377671}.

\bibitem{tao2022}
T.~Tao, \emph{Almost all orbits of the Collatz map attain almost bounded
values}, Forum of Mathematics, Pi \textbf{10} (2022), e12, 1--56.
\href{https://doi.org/10.1017/fmp.2022.8}{doi:10.1017/fmp.2022.8};
\href{https://arxiv.org/abs/1909.03562v7}{arXiv:1909.03562v7}.

\bibitem{lean4322}
Lean Language Reference, \emph{Lean 4.32.2 release notes}, 28 July 2026.
\url{https://lean-lang.org/doc/reference/latest/releases/v4.32.2/}.

\bibitem{lean4340}
Lean Language Reference, \emph{Lean 4.34.0 release notes}, 14 September 2026.
\url{https://lean-lang.org/doc/reference/latest/releases/v4.34.0/}.

\end{thebibliography}
\endgroup

\end{document}
